%==============================================================================
%  BAB II - PERANCANGAN DAN IMPLEMENTASI
%==============================================================================
\section{Perancangan dan Implementasi}

%==============================================================================
\subsection{Perancangan}

\subsubsection{Gambaran Umum Sistem}

Pada Milestone 4, program menambahkan dua tahap baru pada pipeline kompiler:
\emph{Intermediate Code Generation} dan \emph{Execution}. Alur kerjanya dapat
diringkas sebagai berikut.

\begin{center}
\fcolorbox{abuGaris}{abuKode}{\parbox{0.92\textwidth}{\centering\ttfamily
Decorated AST $\rightarrow$ Intermediate Code (TAC) $\rightarrow$ Stack Machine $\rightarrow$ Program Output}}
\end{center}

Generator menerima Decorated AST hasil Milestone 3 lalu menelusurinya untuk
menghasilkan daftar instruksi TAC. Daftar tersebut kemudian dikonversi melalui
sebuah \emph{adapter} dan dieksekusi oleh stack machine, yang mengelola memori
dalam bentuk stack dan mencetak keluaran nyata ketika menemui instruksi
\texttt{WRT}/\texttt{WRTLN}. Spesifikasi prosesnya: \textbf{masukan} berupa
Decorated AST, \textbf{keluaran} berupa Intermediate Code beserta output program,
dengan \textbf{algoritma} penelusuran berbasis DFS untuk pembentukan TAC dan
\emph{stack machine} untuk eksekusi.

\subsubsection{Teknologi yang Digunakan}

Sesuai ketentuan, bahasa pemrograman yang digunakan disamakan dengan milestone
sebelumnya, yaitu \textbf{C++} dengan standar \textbf{C++17}, dikompilasi
menggunakan \texttt{g++} dan diotomasi dengan \textbf{GNU Make}. Tidak ada
pustaka eksternal yang digunakan; seluruh komponen dibangun di atas pustaka
standar C++ (\texttt{<vector>}, \texttt{<string>}, \texttt{<unordered\_map>},
\texttt{<sstream>}, \texttt{<memory>}). Deteksi \emph{overflow} aritmatika
memanfaatkan \emph{builtin} \texttt{\_\_builtin\_*\_overflow} milik \texttt{g++}.

\subsubsection{Struktur Modul dan Arsitektur File}

Sesuai prinsip modularitas, kode Milestone 4 dipisahkan berdasarkan tanggung
jawab ke dalam dua direktori: \texttt{src/intermediate/} (pembentukan TAC) dan
\texttt{src/interpreter/} (eksekusi). Tabel~\ref{tab:file} memerinci berkas-berkas
baru beserta perannya.

\begin{table}[H]
\centering
\caption{Arsitektur berkas Milestone 4.}
\label{tab:file}
\renewcommand{\arraystretch}{1.2}
\begin{tabularx}{\textwidth}{@{}l X@{}}
\toprule
\textbf{Berkas} & \textbf{Tanggung Jawab} \\
\midrule
\texttt{intermediate/tac.\{hpp,cpp\}}     & Kelas \texttt{TacGenerator}: konversi Decorated AST $\rightarrow$ TAC. \\
\texttt{interpreter/runtime\_value.*}      & Struktur \texttt{RuntimeValue}: nilai bertipe pada stack. \\
\texttt{interpreter/runtime\_instruction.hpp} & Struktur \texttt{RuntimeInstruction}: satu instruksi siap eksekusi. \\
\texttt{interpreter/runtime\_context.*}    & Struktur \texttt{RuntimeContext}: stack, IP, BP, dan pembukuan frame. \\
\texttt{interpreter/stack\_machine.*}      & Kelas \texttt{StackMachine}: inti siklus fetch--decode--execute. \\
\texttt{interpreter/runtime\_error\_hook.*}& \texttt{RuntimeErrorHook}: klasifikasi dan pelemparan error runtime. \\
\texttt{interpreter/tac\_adapter.*}        & \texttt{TacAdapter}: jembatan \texttt{TacInstr} $\rightarrow$ \texttt{RuntimeInstruction}. \\
\texttt{interpreter/interpreter.*}         & Kelas \texttt{Interpreter}: pembungkus VM + penangkap error. \\
\bottomrule
\end{tabularx}
\end{table}

\subsubsection{Justifikasi Pilihan Implementasi}

\textbf{Penelusuran berbasis DFS.} Pembentukan TAC dilakukan dengan menelusuri AST
secara \emph{depth-first}. Untuk ekspresi, penelusuran postorder menjamin operand
sudah berada di stack sebelum operatornya dijalankan; untuk pernyataan,
penelusuran preorder yang dipadu pembentukan label menjaga urutan eksekusi sesuai
struktur program.

\textbf{Stack machine, bukan register machine.} Model stack machine dipilih karena
selaras dengan konvensi instruksi yang disediakan (gaya ORKOM) dan menghapus
kebutuhan akan alokasi register/variabel sementara eksplisit---hasil antara cukup
ditinggalkan pada puncak stack.

\textbf{Pola Adapter.} \texttt{TacGenerator} (tim pembentuk kode) dan
\texttt{StackMachine} (tim interpreter) sengaja dipisahkan oleh
\texttt{TacAdapter}. Dengan begitu, kode interpreter tidak bergantung pada detail
internal generator; perubahan pada salah satu sisi tidak merembet ke sisi lain.

\textbf{Error sebagai pengecualian, bukan crash.} Seluruh pelanggaran runtime
dilempar sebagai \texttt{RuntimeError} dan ditangkap di satu titik
(\texttt{Interpreter::run}), sehingga program selalu berhenti secara \emph{graceful}
dengan pesan informatif alih-alih \emph{segmentation fault}.

%==============================================================================
\subsection{Implementasi Intermediate Code Generator}

\texttt{TacGenerator} mengubah Decorated AST menjadi vektor \texttt{TacInstr}.
Tiap instruksi merekam \emph{opcode}, daftar argumen, label opsional, serta target
lompat yang terisi setelah resolusi.

\subsubsection{Pemetaan Alamat Memori}

Sebelum kode dihasilkan, seluruh deklarasi variabel dipindai dan dipetakan ke
alamat memori mulai dari 3 (sel 0--2 dicadangkan untuk static link, dynamic link,
dan return address). Ukuran memori frame utama lalu dialokasikan dengan instruksi
\texttt{INT} sebesar $3 + (\text{jumlah variabel})$.

\begin{lstlisting}[style=kodeCpp, caption={Pemetaan alamat variabel dan alokasi memori awal.}]
void TacGenerator::collectDecls(const AstNodePtr& program) {
    for (auto& child : program->children) {
        if (child->kind != AstKind::Declarations) continue;
        for (auto& decl : child->children) {
            if (decl->kind == AstKind::VarDecl && !decl->value.empty()) {
                string key = lower(decl->value);
                if (addrMap.find(key) == addrMap.end())
                    addrMap[key] = nextAddr++;   // 3, 4, 5, ...
            }
        }
    }
}
// ... pada generate():
int varCount = static_cast<int>(addrMap.size());
emit("INT", {"0", to_string(3 + varCount)});     // INT 0 <3 + jml var>
\end{lstlisting}

\subsubsection{Pembentukan Ekspresi (DFS Postorder)}

Ekspresi dibentuk secara rekursif. Literal menjadi \texttt{LIT}, variabel menjadi
\texttt{LOD} pada alamatnya, sedangkan operasi biner mem-\emph{push} operand kiri
lalu kanan sebelum memanggil \texttt{OPR} dengan opcode yang sesuai. Pemetaan
operator ke opcode mengikuti Tabel~\ref{tab:opr}.

\begin{lstlisting}[style=kodeCpp, caption={Pembentukan ekspresi: hasil ditinggalkan di puncak stack.}]
void TacGenerator::genExpr(const AstNodePtr& node) {
    switch (node->kind) {
        case AstKind::Number: case AstKind::RealNumber:
        case AstKind::BoolLit: /* ... */
            emit("LIT", {"0", node->value});               break;
        case AstKind::Var:
            emit("LOD", {"0", to_string(addrOf(node))});   break;
        case AstKind::BinOp: {
            genExpr(node->children[0]);   // operand kiri di-push lebih dulu
            genExpr(node->children[1]);   // lalu operand kanan
            int code = oprCode(node->value);
            if (code >= 0) emit("OPR", {"0", to_string(code)});
            break;
        }
        case AstKind::UnaryOp:
            genExpr(node->children[0]);
            if (node->value == "-") emit("OPR", {"0", "1"});  // NEG
            break;
    }
}
\end{lstlisting}

\subsubsection{Pembentukan Control Flow}

Pernyataan \texttt{if} dan \texttt{while} dibentuk dengan menyemai label simbolik
(\texttt{L1}, \texttt{L2}, \dots) yang nantinya diresolusi menjadi nomor baris.
Konvensi \emph{IF\_FALSE} diwujudkan oleh \texttt{JPC}: kondisi di-\emph{push},
lalu \texttt{JPC} melompati blok bila kondisi salah.

\begin{lstlisting}[style=kodeCpp, caption={Pembentukan \texttt{if-else} dan \texttt{while} dengan label \& jump.}]
void TacGenerator::genIf(const AstNodePtr& node) {
    string elseLbl = newLabel(), endLbl = newLabel();
    genExpr(cond);                    // push kondisi
    emit("JPC", {"0", elseLbl});      // salah -> lompat ke else
    if (thenNode) genNode(thenNode);
    if (elseNode) emit("JMP", {"0", endLbl});
    emitLabel(elseLbl);
    if (elseNode) genNode(elseNode);
    emitLabel(endLbl);
}

void TacGenerator::genWhile(const AstNodePtr& node) {
    string startLbl = newLabel(), endLbl = newLabel();
    emitLabel(startLbl);
    genExpr(cond);                    // push kondisi
    emit("JPC", {"0", endLbl});       // salah -> keluar loop
    if (body) genNode(body);
    emit("JMP", {"0", startLbl});     // ulangi
    emitLabel(endLbl);
}
\end{lstlisting}

\subsubsection{Resolusi Label dan Pencetakan}

Setelah seluruh kode terbentuk, label diresolusi menjadi nomor baris 1-\emph{based},
lalu \emph{pseudo-op} \texttt{LABEL} dihapus agar daftar instruksi padat dan langsung
dapat dieksekusi (jika dibiarkan, label akan menggeser indeks setiap instruksi
sesudahnya). Saat dicetak, \texttt{JMP}/\texttt{JPC} ditampilkan dengan target
numeriknya dalam format kanonik \texttt{<baris> <OP> <lev> <operand>}.

%==============================================================================
\subsection{Implementasi Interpreter (Stack Machine)}

\subsubsection{Representasi Nilai dan Konteks}

Setiap sel stack menyimpan \texttt{RuntimeValue} yang bertipe (\texttt{INTEGER},
\texttt{REAL}, \texttt{BOOLEAN}, \texttt{CHAR}, \texttt{STRING}, atau
\texttt{NONE}). Tipe yang melekat inilah yang memungkinkan \emph{runtime type
checking} dan pencetakan keluaran yang benar. Seluruh keadaan eksekusi disimpan
dalam \texttt{RuntimeContext}: vektor \texttt{stack}, base pointer \texttt{bp},
return address \texttt{ra}, dynamic link \texttt{dl}, kedalaman frame, instruction
pointer \texttt{ip}, vektor program, serta \emph{buffer} keluaran.

\subsubsection{Siklus Fetch--Decode--Execute}

Metode \texttt{run()} mengulang \texttt{step()} hingga program \emph{halt} atau IP
melewati akhir program. Tiap \texttt{step()} membaca instruksi pada IP lalu
mendispatch ke \emph{handler} sesuai opcode.

\begin{lstlisting}[style=kodeCpp, caption={Inti siklus eksekusi dan dispatch instruksi.}]
void StackMachine::run() {
    while (!context.halted && context.ip < context.program.size())
        step();
}
void StackMachine::execute(const RuntimeInstruction& instr) {
    const string& op = instr.op;
    if (op == "INT") { execINT(instr); return; }
    if (op == "LIT") { execLIT(instr); return; }
    if (op == "LOD") { execLOD(instr); return; }
    if (op == "STO") { execSTO(instr); return; }
    if (op == "JMP") { execJMP(instr); return; }
    if (op == "JPC") { execJPC(instr); return; }
    if (op == "OPR") { execOPR(instr); return; }
    if (op == "RET") { execRET(instr); return; }
    // ... CAL, LABEL ...
    RuntimeErrorHook::invalidOperand(op, "unknown instruction");
}
\end{lstlisting}

\subsubsection{Inisiasi Frame dan Akses Memori}

Instruksi \texttt{INT} membentuk frame baru: mendorong static link, dynamic link,
dan return address, lalu mengalokasikan sel variabel. Akses variabel
(\texttt{LOD}/\texttt{STO}) dihitung relatif terhadap \texttt{bp} dan divalidasi
terhadap batas stack sebelum digunakan.

\begin{lstlisting}[style=kodeCpp, caption={\texttt{LOD} dengan validasi akses memori (mencegah OOB).}]
void StackMachine::execLOD(const RuntimeInstruction& instr) {
    int addr = parseIntArg(instr, 1);
    int absIdx = context.bp + addr;
    if (absIdx < 0 || (size_t)absIdx >= context.stack.size())
        RuntimeErrorHook::invalidMemoryAccess(addr, context.bp,
                                              context.stack.size());
    push(context.stack[absIdx]);
    context.ip++;
}
\end{lstlisting}

\subsubsection{Lompatan dan Konvensi IF\_FALSE}

\texttt{JPC} mengeluarkan kondisi dari puncak stack; bila bernilai \emph{false}, IP
dipindahkan ke target, jika tidak IP maju satu langkah. Setiap target lompat
divalidasi oleh \texttt{resolveJump} agar tidak menunjuk ke luar program.

\begin{lstlisting}[style=kodeCpp, caption={\texttt{JPC} (jump-if-false) dan validasi target lompat.}]
void StackMachine::execJPC(const RuntimeInstruction& instr) {
    RuntimeValue cond = pop();
    if (!cond.isTruthy()) context.ip = resolveJump(instr.target);
    else                  context.ip++;
}
size_t StackMachine::resolveJump(int target) const {
    int n = (int)context.program.size();
    if (target < 1 || target > n)
        RuntimeErrorHook::invalidJumpTarget(target, n);
    return (size_t)(target - 1);   // 1-based -> 0-based
}
\end{lstlisting}

\subsubsection{Operasi melalui OPR}

\texttt{execOPR} mendispatch nomor operasi ke sub-handler. Operasi aritmatika
mengeluarkan dua operand (atau satu untuk \texttt{NEG}), memvalidasi tipenya,
melakukan perhitungan ber-\emph{checked} terhadap overflow, lalu mem-\emph{push}
hasilnya. Operasi perbandingan menghasilkan \texttt{BOOLEAN}, sedangkan
\texttt{WRT}/\texttt{WRTLN} menulis nilai puncak stack ke \emph{buffer} keluaran.

%==============================================================================
\subsection{Implementasi Error Handling dan Keamanan Runtime}

Seluruh kondisi error dipusatkan pada modul \texttt{RuntimeErrorHook}. Setiap
pelanggaran memanggil fungsi yang memformat pesan \texttt{[ERROR] <Jenis>: <detail> !}
lalu melempar \texttt{RuntimeError} dengan klasifikasi jenisnya, sehingga pesan
konsisten dan mudah ditelusuri.

\textbf{Aritmatika ber-checked.} Operasi integer memakai \emph{builtin} g++ untuk
mendeteksi \emph{wrap-around}, lalu melempar \texttt{OverflowError}/\texttt{UnderflowError}.

\begin{lstlisting}[style=kodeCpp, caption={Penjumlahan dengan deteksi overflow/underflow.}]
long long StackMachine::addChecked(long long a, long long b, const string& op) {
    long long r;
    if (__builtin_add_overflow(a, b, &r)) {
        if (a > 0 && b > 0) RuntimeErrorHook::numericOverflow(op);
        RuntimeErrorHook::numericUnderflow(op);
    }
    return r;
}
\end{lstlisting}

\textbf{Deteksi stack corruption dan smashing.} Saat \texttt{INT}, sebuah
\texttt{FrameGuard} mencatat posisi puncak yang seharusnya saat keluar
(\emph{expected top}) dan salinan return address sebagai \emph{canary}. Pada
\texttt{RET}, posisi puncak aktual dibandingkan dengan \emph{expected top} (deteksi
\emph{corruption}) dan sel return address dibandingkan dengan \emph{canary}
(deteksi \emph{smashing}).

\begin{lstlisting}[style=kodeCpp, caption={Verifikasi keseimbangan frame saat \texttt{RET}.}]
const FrameGuard guard = context.frames.back();
context.frames.pop_back();
int actualTop = (int)context.stack.size();
if (actualTop != guard.expectedTop)              // push/pop tak simetris
    RuntimeErrorHook::stackCorruption("unbalanced push/pop ...",
                                      guard.expectedTop, actualTop);
int raIdx = guard.basePtr + 2;                    // sel return address
if (context.stack[raIdx].intVal != guard.canaryRa)
    RuntimeErrorHook::stackSmashing(guard.canaryRa, context.stack[raIdx].intVal);
\end{lstlisting}

Tabel~\ref{tab:err} merangkum jenis-jenis error runtime yang ditangani beserta
mekanisme deteksinya.

\begin{table}[H]
\centering
\caption{Ringkasan penanganan error runtime.}
\label{tab:err}
\renewcommand{\arraystretch}{1.15}
\begin{tabularx}{\textwidth}{@{}l X@{}}
\toprule
\textbf{Jenis Error} & \textbf{Mekanisme Deteksi} \\
\midrule
DivisionByZeroError      & Pembagi bernilai nol pada \texttt{DIV}/\texttt{MOD}. \\
OverflowError / UnderflowError & Hasil aritmatika integer membungkus (\emph{wrap-around}). \\
StackOverflowError       & Kedalaman frame atau ukuran stack melampaui batas. \\
StackUnderflowError      & \emph{Pop} dipanggil saat stack kekurangan nilai. \\
InvalidJumpError         & Target \texttt{JMP}/\texttt{JPC}/\texttt{CAL} di luar program. \\
InvalidMemoryAccessError & Alamat \texttt{LOD}/\texttt{STO} di luar batas stack. \\
StackCorruptionError     & Posisi stack pointer tidak sesuai saat keluar frame. \\
StackSmashingError       & \emph{Canary} return address tertimpa data lain. \\
TypeMismatchError        & Operand aritmatika ternyata bukan nilai numerik. \\
InvalidOperandError      & Argumen instruksi hilang atau bukan integer valid. \\
\bottomrule
\end{tabularx}
\end{table}
